\documentclass{article}
\usepackage[T1]{fontenc}
\usepackage{lmodern}
\usepackage{graphicx}
\usepackage{amsmath,amssymb}
\usepackage[a4paper,margin=2cm]{geometry}
\usepackage[absolute,overlay]{textpos}
\setlength{\TPHorizModule}{1mm}
\setlength{\TPVertModule}{1mm}
\usepackage[indonesian]{babel}
\usepackage{hyperref}
\setlength{\emergencystretch}{2em}
% unit: Halaman 71

\begin{document}

% === Halaman 71 (llm) ===

\section*{Keamanan A5}

\begin{itemize}
    \item Keamanan A5/1 terletak pada pembangkitan bit-bit acak yang dihasilkan oleh tiga buah LFSR di atas.
    \item Tujuan kriptanalisis A5/1 adalah untuk menemukan kunci sesi yang digunakan dalam pembangkitan bit-bit acak.
    \item Dalam serangan tersebut diasumsikan penyerang mengetahui luaran A5/1 pada periode awal komunikasi, untuk selanjutnya menemukan kunci sesi yang digunakan untuk mendekripsi sisa komunikasi selanjutnya.
\end{itemize}

\end{document}
